\section{Introduction}
An exchange can help its recipient and create a dangerous competitor. Selling the final deed in a Monopoly color group supplies cash today while enabling construction against the seller tomorrow. Refusing a sale preserves a blocker, but the holder may later lose it to someone else through bankruptcy. The same dependence appears within an agent's own finances. Selling buildings sacrifices income yet makes previously unavailable mortgages legal. These decisions concern how an action changes subsequent opportunities.

Such settings allow a useful separation between what an agent says, what its state permits, and what it actually does. A statement can correctly identify an opportunity without leading to a profitable choice. An agent can reach a financial goal while receiving substantial intermediate feedback. Conversely, an incorrect explanation can accompany a favorable transaction. Studying these objects together supplies a more precise account of economic behavior than judging a plausible explanation or a final win alone.

We examine twenty complete MonopolyBench games played by GPT-6.1 Sol, Gemini 4 Argon, Claude Opus 5.5 and Grok 4.7. The corpus records 6,989 decisions and the resulting 50,918 events. The study follows bargaining into financing, construction, rent and ownership transfer. Every offer has a recipient-facing state and response; each liquidation episode has a specified obligation and a payment or bankruptcy endpoint. Persistent consequences make the games informative even when no defensible value reference exists for a rejected trade.

The clearest repeated positive finding concerns debts that cannot be paid using cash, all building salvage and the mortgages listed at entry. Sixteen encountered episodes require additional collateral released by selling buildings. Fourteen reach payment. The successful entry texts mostly describe that dependency before the menu changes; the two failures omit the released collateral and declare bankruptcy immediately. This is evidence of prospective financial assessment followed by dependent execution. It is more specific than a general claim that the agents plan, and more informative than a count of successful payments without their required resources.

Bargaining shows a related concern with other players' opportunities, but its aggregate interpretation is less settled. Opening swaps completing only the proposer are rejected more often than swaps completing neither party. The contrast depends materially on who proposes the offers. Pure cash bids do not have the same observed ordering. We present the complete form and completion table, proposer and recipient strata, and sensitivity across the five seed blocks. The analysis describes an exploratory association among unequal offers; it does not isolate the effect of enabling a rival's development.

Three main incidents clarify the connection between economic statements and action. A player refuses its own complete group while citing the rival's building cash. A seller presents a purchase as protection against a third party's inheritance while omitting available collateral. A debtor announces an aggregate liquidation budget, executes ten actions and retains the predicted \$14. Five further incidents preserve pricing, a public/private bargaining ceiling, factual correction, attempted cooperation and withdrawal of an assurance. They supply distinct observations, not a frequency estimate of deceptive conduct.

The contribution is an observational account of economic opportunity assessment at two evidential levels. Local debt has a verifiable immediate target, enabling statements and actions to be checked against resource arithmetic. Bargaining exposes expressed rival concern and actual response patterns while leaving the value of each choice unresolved. We retain that asymmetry throughout. The broader motivating questions concern reasoning, opponent modeling, negotiation and deception, planning and memory, adaptation, risk, collective dynamics and evaluation validity. Appendix~\ref{app:questions} answers their measured components, together with questions about talk versus action, general capability and skill versus chance.
